\documentclass[10pt,a4paper]{amsart}
\usepackage{amsmath,amssymb,amsthm,graphicx}
\usepackage{verbatim}
\newenvironment{reference}{\comment}{\endcomment}
%\newenvironment{reference}{}{}
\newenvironment{slogan}{\comment}{\endcomment}
\newenvironment{history}{\comment}{\endcomment}

% For commutative diagrams we use Xy-pic
\usepackage[all]{xy}

% We use 2cell for 2-commutative diagrams.
\xyoption{2cell}
% Theorem environments.
%
\theoremstyle{plain}
\newtheorem{theorem}[subsection]{Theorem}
\newtheorem{proposition}[subsection]{Proposition}
\newtheorem{lemma}[subsection]{Lemma}

\theoremstyle{definition}
\newtheorem{definition}[subsection]{Definition}
\newtheorem{example}[subsection]{Example}
\newtheorem{exercise}[subsection]{Exercise}
\newtheorem{situation}[subsection]{Situation}

\theoremstyle{remark}
\newtheorem{remark}[subsection]{Remark}
\newtheorem{remarks}[subsection]{Remarks}

\numberwithin{equation}{subsection}

% Macros
%
\def\lim{\mathop{\mathrm{lim}}\nolimits}
\def\colim{\mathop{\mathrm{colim}}\nolimits}
\def\Spec{\mathop{\mathrm{Spec}}}
\def\Hom{\mathop{\mathrm{Hom}}\nolimits}
\def\Ext{\mathop{\mathrm{Ext}}\nolimits}
\def\SheafHom{\mathop{\mathcal{H}\!\mathit{om}}\nolimits}
\def\SheafExt{\mathop{\mathcal{E}\!\mathit{xt}}\nolimits}
\def\Sch{\mathit{Sch}}
\def\Mor{\mathop{\mathrm{Mor}}\nolimits}
\def\Ob{\mathop{\mathrm{Ob}}\nolimits}
\def\Sh{\mathop{\mathit{Sh}}\nolimits}
\def\NL{\mathop{N\!L}\nolimits}
\def\CH{\mathop{\mathrm{CH}}\nolimits}
\def\proetale{{pro\text{-}\acute{e}tale}}
\def\etale{{\acute{e}tale}}
\def\QCoh{\mathit{QCoh}}
\def\Ker{\mathop{\mathrm{Ker}}}
\def\Im{\mathop{\mathrm{Im}}}
\def\Coker{\mathop{\mathrm{Coker}}}
\def\Coim{\mathop{\mathrm{Coim}}}

% Boxtimes
%
\DeclareMathSymbol{\boxtimes}{\mathbin}{AMSa}{"02}

%
% Macros for moduli stacks/spaces
%
\def\QCohstack{\mathcal{QC}\!\mathit{oh}}
\def\Cohstack{\mathcal{C}\!\mathit{oh}}
\def\Spacesstack{\mathcal{S}\!\mathit{paces}}
\def\Quotfunctor{\mathrm{Quot}}
\def\Hilbfunctor{\mathrm{Hilb}}
\def\Curvesstack{\mathcal{C}\!\mathit{urves}}
\def\Polarizedstack{\mathcal{P}\!\mathit{olarized}}
\def\Complexesstack{\mathcal{C}\!\mathit{omplexes}}
% \Pic is the operator that assigns to X its picard group, usage \Pic(X)
% \Picardstack_{X/B} denotes the Picard stack of X over B
% \Picardfunctor_{X/B} denotes the Picard functor of X over B
\def\Pic{\mathop{\mathrm{Pic}}\nolimits}
\def\Picardstack{\mathcal{P}\!\mathit{ic}}
\def\Picardfunctor{\mathrm{Pic}}
\def\Deformationcategory{\mathcal{D}\!\mathit{ef}}
\usepackage{cleveref}
\crefname{originalreference}{original reference}{original references}
\Crefname{originalreference}{Original reference}{Original references}
\makeatletter
\newlabel{definition-big-proetale-site}{{098G}{source}{}{sourcepointer-1}{}}
\newlabel{definition-big-proetale-site@cref}{{[originalreference][0][]098G}{[1][0][]source}{}{}{}}
\makeatother
\crefname{originalreference}{original reference}{original references}
\Crefname{originalreference}{Original reference}{Original references}
\makeatletter
\newlabel{lemma-get-many-weakly-contractible}{{098I}{source}{}{sourcepointer-2}{}}
\newlabel{lemma-get-many-weakly-contractible@cref}{{[originalreference][0][]098I}{[1][0][]source}{}{}{}}
\makeatother
\crefname{originalreference}{original reference}{original references}
\Crefname{originalreference}{Original reference}{Original references}
\makeatletter
\newlabel{lemma-w-contractible-proetale-cover}{{098F}{source}{}{sourcepointer-3}{}}
\newlabel{lemma-w-contractible-proetale-cover@cref}{{[originalreference][0][]098F}{[1][0][]source}{}{}{}}
\makeatother
\title{Every pro-etale covering of an object of the big pro-etale site is refined by a covering of the site}
\author{The Stacks Project Authors}
\date{stacks-project:proetale.tex:local-snapshot-20261012; GNU FDL 1.2 or later}
\begin{document}
\maketitle
\noindent\textit{Editorial note.} All delivered mathematical text is the author's own, copied from the pinned Stacks Project chapter file proetale.tex; no statement or proof step was rewritten, reordered or removed.\par\medskip
\noindent\textit{Editorial note.} The counted proof is the author's complete argument: the covering by spectra of w-contractible rings, the refinement of each induced pro-etale covering by a finite disjoint open covering, and the collection of the pieces into a covering of the site.\par\medskip
\noindent\textit{Editorial note.} The label lemma-proetale-induced names the counted statement, and the proof environment immediately following it in the source is delivered verbatim as its proof.\par\medskip
\noindent\textit{Editorial note.} The counted statement is carried by the author's own lemma environment in the chapter source.\par\medskip
\noindent\textit{Editorial note.} Pre-open gate: python operations/agent-authoring/tools/precard\_gate.py --source proetale.tex --label lemma-proetale-induced --start 6515 --end 6534 --local-id 28007 returned RESULT=OK (L28007 free, occupancy FREE, public interval clear). Candidate taken from sources/stacks-proetale/candidates-halfB.json.\par\medskip
\noindent\textit{Editorial note.} The statement and proof are copied from the pinned chapter file; the cited inputs are kept as external pointers with their canonical Stacks tags.\par\medskip
\noindent\textit{Allowed standard prerequisites (editorial).}\begin{itemize}
\item Coverings by spectra of w-contractible rings. Same chapter, the lemma producing a covering of any object by affine objects which are spectra of w-contractible rings. Supplies the covering \{V\_k -> T\} whose members are weakly contractible.
\item Refinement of pro-etale coverings of a w-contractible affine scheme by finite disjoint open coverings. Same chapter, the lemma that every pro-etale covering of a w-contractible affine scheme is refined by a finite Zariski covering by a disjoint union. Supplies the refinement of each pulled back covering of V\_k by a finite disjoint open covering.
\end{itemize}\medskip
\ifdefined\hypertarget\hypertarget{sourcepointer-1}{}\fi
\noindent\textit{Source reference.} Original reference 098G: sources/stacks-proetale/pinned-original-source.tex, source lines 2465--2465. Prerequisite pointer invoked at proof line(s) 6518; it names a standard fact of the same chapter that the counted argument uses, and every step of the counted proof built on it is retained verbatim in the delivered body.\par\medskip
\ifdefined\hypertarget\hypertarget{sourcepointer-2}{}\fi
\noindent\textit{Source reference.} Original reference 098I: sources/stacks-proetale/pinned-original-source.tex, source lines 3087--3087. Prerequisite pointer invoked at proof line(s) 6527; it names a standard fact of the same chapter that the counted argument uses, and every step of the counted proof built on it is retained verbatim in the delivered body.\par\medskip
\ifdefined\hypertarget\hypertarget{sourcepointer-3}{}\fi
\noindent\textit{Source reference.} Original reference 098F: sources/stacks-proetale/pinned-original-source.tex, source lines 2998--2998. Prerequisite pointer invoked at proof line(s) 6531; it names a standard fact of the same chapter that the counted argument uses, and every step of the counted proof built on it is retained verbatim in the delivered body.\par\medskip
\begin{lemma}
\label{lemma-proetale-induced}
Let $\Sch_\proetale$ be a big pro-\'etale site as in
Definition \ref{definition-big-proetale-site}.
Let $T \in \Ob(\Sch_\proetale)$.
Let $\{T_i \to T\}_{i \in I}$ be an arbitrary pro-\'etale covering of $T$.
There exists a covering $\{U_j \to T\}_{j \in J}$ of $T$ in the site
$\Sch_\proetale$ which refines $\{T_i \to T\}_{i \in I}$.
\end{lemma}

\begin{proof}
Namely, we first let $\{V_k \to T\}$ be a covering as in
Lemma \ref{lemma-get-many-weakly-contractible}.
Then the pro-\'etale coverings $\{T_i \times_T V_k \to V_k\}$
can be refined by a finite disjoint open covering
$V_k = V_{k, 1} \amalg \ldots \amalg V_{k, n_k}$, see
Lemma \ref{lemma-w-contractible-proetale-cover}.
Then $\{V_{k, i} \to T\}$ is a covering of $\Sch_\proetale$
which refines $\{T_i \to T\}_{i \in I}$.
\end{proof}

\end{document}
